% Options for packages loaded elsewhere
% Options for packages loaded elsewhere
\PassOptionsToPackage{unicode}{hyperref}
\PassOptionsToPackage{hyphens}{url}
\PassOptionsToPackage{dvipsnames,svgnames,x11names}{xcolor}
%
\documentclass[
  russian,
  letterpaper,
  DIV=11,
  numbers=noendperiod]{scrartcl}
\usepackage{xcolor}
\usepackage{amsmath,amssymb}
\setcounter{secnumdepth}{-\maxdimen} % remove section numbering
\usepackage{iftex}
\ifPDFTeX
  \usepackage[T1]{fontenc}
  \usepackage[utf8]{inputenc}
  \usepackage{textcomp} % provide euro and other symbols
\else % if luatex or xetex
  \usepackage{unicode-math} % this also loads fontspec
  \defaultfontfeatures{Scale=MatchLowercase}
  \defaultfontfeatures[\rmfamily]{Ligatures=TeX,Scale=1}
\fi
\usepackage{lmodern}
\ifPDFTeX\else
  % xetex/luatex font selection
  \setmainfont[]{Liberation Serif}
  \setsansfont[]{Liberation Sans}
  \setmonofont[]{Liberation Mono}
\fi
% Use upquote if available, for straight quotes in verbatim environments
\IfFileExists{upquote.sty}{\usepackage{upquote}}{}
\IfFileExists{microtype.sty}{% use microtype if available
  \usepackage[]{microtype}
  \UseMicrotypeSet[protrusion]{basicmath} % disable protrusion for tt fonts
}{}
\makeatletter
\@ifundefined{KOMAClassName}{% if non-KOMA class
  \IfFileExists{parskip.sty}{%
    \usepackage{parskip}
  }{% else
    \setlength{\parindent}{0pt}
    \setlength{\parskip}{6pt plus 2pt minus 1pt}}
}{% if KOMA class
  \KOMAoptions{parskip=half}}
\makeatother
% Make \paragraph and \subparagraph free-standing
\makeatletter
\ifx\paragraph\undefined\else
  \let\oldparagraph\paragraph
  \renewcommand{\paragraph}{
    \@ifstar
      \xxxParagraphStar
      \xxxParagraphNoStar
  }
  \newcommand{\xxxParagraphStar}[1]{\oldparagraph*{#1}\mbox{}}
  \newcommand{\xxxParagraphNoStar}[1]{\oldparagraph{#1}\mbox{}}
\fi
\ifx\subparagraph\undefined\else
  \let\oldsubparagraph\subparagraph
  \renewcommand{\subparagraph}{
    \@ifstar
      \xxxSubParagraphStar
      \xxxSubParagraphNoStar
  }
  \newcommand{\xxxSubParagraphStar}[1]{\oldsubparagraph*{#1}\mbox{}}
  \newcommand{\xxxSubParagraphNoStar}[1]{\oldsubparagraph{#1}\mbox{}}
\fi
\makeatother


\usepackage{longtable,booktabs,array}
\usepackage{calc} % for calculating minipage widths
% Correct order of tables after \paragraph or \subparagraph
\usepackage{etoolbox}
\makeatletter
\patchcmd\longtable{\par}{\if@noskipsec\mbox{}\fi\par}{}{}
\makeatother
% Allow footnotes in longtable head/foot
\IfFileExists{footnotehyper.sty}{\usepackage{footnotehyper}}{\usepackage{footnote}}
\makesavenoteenv{longtable}
\usepackage{graphicx}
\makeatletter
\newsavebox\pandoc@box
\newcommand*\pandocbounded[1]{% scales image to fit in text height/width
  \sbox\pandoc@box{#1}%
  \Gscale@div\@tempa{\textheight}{\dimexpr\ht\pandoc@box+\dp\pandoc@box\relax}%
  \Gscale@div\@tempb{\linewidth}{\wd\pandoc@box}%
  \ifdim\@tempb\p@<\@tempa\p@\let\@tempa\@tempb\fi% select the smaller of both
  \ifdim\@tempa\p@<\p@\scalebox{\@tempa}{\usebox\pandoc@box}%
  \else\usebox{\pandoc@box}%
  \fi%
}
% Set default figure placement to htbp
\def\fps@figure{htbp}
\makeatother



\ifLuaTeX
\usepackage[bidi=basic,provide=*]{babel}
\else
\usepackage[bidi=default,provide=*]{babel}
\fi
\ifPDFTeX
\else
\babelfont{rm}[]{Liberation Serif}
\fi
% get rid of language-specific shorthands (see #6817):
\let\LanguageShortHands\languageshorthands
\def\languageshorthands#1{}


\setlength{\emergencystretch}{3em} % prevent overfull lines

\providecommand{\tightlist}{%
  \setlength{\itemsep}{0pt}\setlength{\parskip}{0pt}}



 


\usepackage{fontspec}
\setmainfont{Liberation Serif}
\setsansfont{Liberation Sans}
\setmonofont{Liberation Mono}
\KOMAoption{captions}{tableheading}
\makeatletter
\@ifpackageloaded{caption}{}{\usepackage{caption}}
\AtBeginDocument{%
\ifdefined\contentsname
  \renewcommand*\contentsname{Содержание}
\else
  \newcommand\contentsname{Содержание}
\fi
\ifdefined\listfigurename
  \renewcommand*\listfigurename{Список иллюстраций}
\else
  \newcommand\listfigurename{Список иллюстраций}
\fi
\ifdefined\listtablename
  \renewcommand*\listtablename{Список таблиц}
\else
  \newcommand\listtablename{Список таблиц}
\fi
\ifdefined\figurename
  \renewcommand*\figurename{Рисунок}
\else
  \newcommand\figurename{Рисунок}
\fi
\ifdefined\tablename
  \renewcommand*\tablename{Таблица}
\else
  \newcommand\tablename{Таблица}
\fi
}
\@ifpackageloaded{float}{}{\usepackage{float}}
\floatstyle{ruled}
\@ifundefined{c@chapter}{\newfloat{codelisting}{h}{lop}}{\newfloat{codelisting}{h}{lop}[chapter]}
\floatname{codelisting}{Список}
\newcommand*\listoflistings{\listof{codelisting}{Листинги}}
\makeatother
\makeatletter
\makeatother
\makeatletter
\@ifpackageloaded{caption}{}{\usepackage{caption}}
\@ifpackageloaded{subcaption}{}{\usepackage{subcaption}}
\makeatother
\usepackage{bookmark}
\IfFileExists{xurl.sty}{\usepackage{xurl}}{} % add URL line breaks if available
\urlstyle{same}
\hypersetup{
  pdftitle={Лабораторная работа № 1},
  pdfauthor={Левищева Анастасия Петровна},
  pdflang={ru},
  colorlinks=true,
  linkcolor={blue},
  filecolor={Maroon},
  citecolor={Blue},
  urlcolor={Blue},
  pdfcreator={LaTeX via pandoc}}


\title{Лабораторная работа № 1}
\usepackage{etoolbox}
\makeatletter
\providecommand{\subtitle}[1]{% add subtitle to \maketitle
  \apptocmd{\@title}{\par {\large #1 \par}}{}{}
}
\makeatother
\subtitle{Архитектура компьютеров и операционные системы}
\author{Левищева Анастасия Петровна}
\date{2026-03-07}
\begin{document}
\maketitle


\section{Информация}\label{ux438ux43dux444ux43eux440ux43cux430ux446ux438ux44f}

\subsection{Докладчик}\label{ux434ux43eux43aux43bux430ux434ux447ux438ux43a}

\begin{itemize}
\tightlist
\item
  Левищева Анастасия Петровна
\item
  group NKAbd-02-25
\item
  Российский университет дружбы народов им. П. Лумумбы
\item
  {[}1032252643@rudn.ru{]}
\end{itemize}

\section{Цель
работы}\label{ux446ux435ux43bux44c-ux440ux430ux431ux43eux442ux44b}

Целью данной работы является приобретение практических навыков установки
операционной системы на виртуальную машину, настройки минимально
необходимых для дальнейшей работы сервисов.

\section{Задание}\label{ux437ux430ux434ux430ux43dux438ux435}

\begin{verbatim}
Дождитесь загрузки графического окружения и откройте терминал. В окне терминала проанализируйте последовательность загрузки системы, выполнив команду dmesg. Можно просто просмотреть вывод этой команды:

dmesg | less

Можно использовать поиск с помощью grep:

dmesg | grep -i "то, что ищем"

Получите следующую информацию.
    Версия ядра Linux (Linux version).
    Частота процессора (Detected Mhz processor).
    Модель процессора (CPU0).
    Объём доступной оперативной памяти (Memory available).
    Тип обнаруженного гипервизора (Hypervisor detected).
    Тип файловой системы корневого раздела.
    Последовательность монтирования файловых систем.
\end{verbatim}

\section{Выполнение лабораторной
работы}\label{ux432ux44bux43fux43eux43bux43dux435ux43dux438ux435-ux43bux430ux431ux43eux440ux430ux442ux43eux440ux43dux43eux439-ux440ux430ux431ux43eux442ux44b}

\subsection{Устанавливаем виртуальную
машину.}\label{ux443ux441ux442ux430ux43dux430ux432ux43bux438ux432ux430ux435ux43c-ux432ux438ux440ux442ux443ux430ux43bux44cux43dux443ux44e-ux43cux430ux448ux438ux43dux443.}

Папка виртуальных машин. Значения по умолчанию Linux: \$HOME/VirtualBox
VMs.

\subsection{Установочный
образ}\label{ux443ux441ux442ux430ux43dux43eux432ux43eux447ux43dux44bux439-ux43eux431ux440ux430ux437}

Перенесем установочный образ в папку /var/tmp/имя\_пользователя/iso:
mkdir -p ``\$(vboxmanage list systemproperties \textbar{} grep `Default
machine folder:' \textbar{} cut -d':' -f2 \textbar{} tr -d ' ')/iso''

mv Fedora-Sway-Live-x86\_64-39-1.5.iso ``\$(vboxmanage list
systemproperties \textbar{} grep `Default machine folder:' \textbar{}
cut -d':' -f2 \textbar{} tr -d ' ')/iso''

\subsection{Настройка
хост-клавиши}\label{ux43dux430ux441ux442ux440ux43eux439ux43aux430-ux445ux43eux441ux442-ux43aux43bux430ux432ux438ux448ux438}

Хост-клавишей по умолчанию является правый Ctrl.

По умолчанию в дисплейных классах на клавише правый Ctrl находится
переключатель языка ввода.

Эти значения могут конфликтовать.

\subsection{Графический
интерфейс}\label{ux433ux440ux430ux444ux438ux447ux435ux441ux43aux438ux439-ux438ux43dux442ux435ux440ux444ux435ux439ux441}

В меню выберите Файл, Настройки.

Выберите Ввод, вкладка Виртуальная машина.

Выберите Сочетание клавиш в строке Хост-комбинация.

Нажмите новое сочетание клавиш.

Нажмите ОК, чтобы сохранить изменения.

\subsection{Командная
строка}\label{ux43aux43eux43cux430ux43dux434ux43dux430ux44f-ux441ux442ux440ux43eux43aux430}

Проверим текущую комбинацию для хост-клавиши:

VBoxManage getextradata global GUI/Input/HostKeyCombination

По умолчанию установлена комбинация 65508, соответствующая правой
клавише Ctrl.

Установите нужную клавишу (в примере клавиша Menu): VBoxManage
setextradata global GUI/Input/HostKeyCombination 65383

\subsubsection{\texorpdfstring{Создадим и запустим виртуальную машину.
Запускаем установочный образ:
(рис.~\ref{fig-001}).}{Создадим и запустим виртуальную машину. Запускаем установочный образ: (рис.~).}}\label{ux441ux43eux437ux434ux430ux434ux438ux43c-ux438-ux437ux430ux43fux443ux441ux442ux438ux43c-ux432ux438ux440ux442ux443ux430ux43bux44cux43dux443ux44e-ux43cux430ux448ux438ux43dux443.-ux437ux430ux43fux443ux441ux43aux430ux435ux43c-ux443ux441ux442ux430ux43dux43eux432ux43eux447ux43dux44bux439-ux43eux431ux440ux430ux437-ux440ux438ux441.-fig-001.}

\begin{figure}

\centering{

\includegraphics[width=0.7\linewidth,height=\textheight,keepaspectratio]{image/1.jpg}

}

\caption{\label{fig-001}Запуск установочного образа}

\end{figure}%

\subsubsection{\texorpdfstring{После установки:
(рис.~\ref{fig-002}).}{После установки: (рис.~).}}\label{ux43fux43eux441ux43bux435-ux443ux441ux442ux430ux43dux43eux432ux43aux438-ux440ux438ux441.-fig-002.}

\begin{figure}

\centering{

\includegraphics[width=0.7\linewidth,height=\textheight,keepaspectratio]{image/2.jpg}

}

\caption{\label{fig-002}Вход в ОС под заданной при установке учетной
записью.}

\end{figure}%

\subsubsection{Войдите в ОС под заданной вами при установке учётной
записью.}\label{ux432ux43eux439ux434ux438ux442ux435-ux432-ux43eux441-ux43fux43eux434-ux437ux430ux434ux430ux43dux43dux43eux439-ux432ux430ux43cux438-ux43fux440ux438-ux443ux441ux442ux430ux43dux43eux432ux43aux435-ux443ux447ux451ux442ux43dux43eux439-ux437ux430ux43fux438ux441ux44cux44e.}

Нажмите комбинацию Win+Enter для запуска терминала.

Переключитесь на роль супер-пользователя: sudo -i

\subsection{Обновления}\label{ux43eux431ux43dux43eux432ux43bux435ux43dux438ux44f}

\subsubsection{Установите средства
разработки:}\label{ux443ux441ux442ux430ux43dux43eux432ux438ux442ux435-ux441ux440ux435ux434ux441ux442ux432ux430-ux440ux430ux437ux440ux430ux431ux43eux442ux43aux438}

sudo dnf -y group install development-tools (рис.~\ref{fig-003}).

\begin{figure}

\centering{

\includegraphics[width=0.7\linewidth,height=\textheight,keepaspectratio]{image/3.jpg}

}

\caption{\label{fig-003}Установка средств разработки}

\end{figure}%

\subsubsection{Обновить все
пакеты}\label{ux43eux431ux43dux43eux432ux438ux442ux44c-ux432ux441ux435-ux43fux430ux43aux435ux442ux44b}

sudo dnf -y update

\subsection{Повышение комфорта
работы}\label{ux43fux43eux432ux44bux448ux435ux43dux438ux435-ux43aux43eux43cux444ux43eux440ux442ux430-ux440ux430ux431ux43eux442ux44b}

Программы для удобства работы в консоли: sudo dnf -y install tmux mc

Другой вариант консоли: sudo dnf -y install kitty

\subsection{Автоматическое
обновление}\label{ux430ux432ux442ux43eux43cux430ux442ux438ux447ux435ux441ux43aux43eux435-ux43eux431ux43dux43eux432ux43bux435ux43dux438ux435}

При необходимости можно использовать автоматическое обновление.

Установка программного обеспечения: sudo dnf -y install dnf-automatic

Задаёте необходимую конфигурацию в файле /etc/dnf/automatic.conf.

Запустите таймер: sudo systemctl enable --now dnf-automatic.timer

\subsection{Отключение
SELinux}\label{ux43eux442ux43aux43bux44eux447ux435ux43dux438ux435-selinux}

В данном курсе мы не будем рассматривать работу с системой безопасности
SELinux.

Поэтому отключим его.

\subsubsection{В файле /etc/selinux/config замените
значение}\label{ux432-ux444ux430ux439ux43bux435-etcselinuxconfig-ux437ux430ux43cux435ux43dux438ux442ux435-ux437ux43dux430ux447ux435ux43dux438ux435}

SELINUX=enforcing

на значение

SELINUX=permissive (рис.~\ref{fig-004}).

\begin{figure}

\centering{

\includegraphics[width=0.7\linewidth,height=\textheight,keepaspectratio]{image/4.jpg}

}

\caption{\label{fig-004}Отключение безопасности SELinux}

\end{figure}%

\subsubsection{Перегрузите виртуальную
машину:}\label{ux43fux435ux440ux435ux433ux440ux443ux437ux438ux442ux435-ux432ux438ux440ux442ux443ux430ux43bux44cux43dux443ux44e-ux43cux430ux448ux438ux43dux443}

sudo systemctl reboot

\subsection{Настройка раскладки
клавиатуры}\label{ux43dux430ux441ux442ux440ux43eux439ux43aux430-ux440ux430ux441ux43aux43bux430ux434ux43aux438-ux43aux43bux430ux432ux438ux430ux442ux443ux440ux44b}

Войдите в ОС под заданной вами при установке учётной записью.

Нажмите комбинацию Win+Enter для запуска терминала.

Запустите терминальный мультиплексор tmux: tmux

Создайте конфигурационный файл
\textasciitilde/.config/sway/config.d/95-system-keyboard-config.conf:
mkdir -p \textasciitilde/.config/sway touch
\textasciitilde/.config/sway/config.d/95-system-keyboard-config.conf

\subsubsection{Отредактируйте конфигурационный
файл}\label{ux43eux442ux440ux435ux434ux430ux43aux442ux438ux440ux443ux439ux442ux435-ux43aux43eux43dux444ux438ux433ux443ux440ux430ux446ux438ux43eux43dux43dux44bux439-ux444ux430ux439ux43b}

\textasciitilde/.config/sway/config.d/95-system-keyboard-config.conf:
exec\_always /usr/libexec/sway-systemd/locale1-xkb-config --oneshot
(рис.~\ref{fig-005}).

\begin{figure}

\centering{

\includegraphics[width=0.7\linewidth,height=\textheight,keepaspectratio]{image/5.jpg}

}

\caption{\label{fig-005}Конфигурационный файл}

\end{figure}%

\subsubsection{Переключитесь на роль
супер-пользователя:}\label{ux43fux435ux440ux435ux43aux43bux44eux447ux438ux442ux435ux441ux44c-ux43dux430-ux440ux43eux43bux44c-ux441ux443ux43fux435ux440-ux43fux43eux43bux44cux437ux43eux432ux430ux442ux435ux43bux44f}

sudo -i

Отредактируйте конфигурационный файл
/etc/X11/xorg.conf.d/00-keyboard.conf: Section ``InputClass'' Identifier
``system-keyboard'' MatchIsKeyboard ``on'' Option ``XkbLayout''
``us,ru'' Option ``XkbVariant'' ``,winkeys'' Option ``XkbOptions''
``grp:rctrl\_toggle,compose:ralt,terminate:ctrl\_alt\_bksp'' EndSection

Для этого можно использовать файловый менеджер mc и его встроенный
редактор.

Перегрузите виртуальную машину: sudo systemctl reboot

\subsection{Установка имени пользователя и названия
хоста}\label{ux443ux441ux442ux430ux43dux43eux432ux43aux430-ux438ux43cux435ux43dux438-ux43fux43eux43bux44cux437ux43eux432ux430ux442ux435ux43bux44f-ux438-ux43dux430ux437ux432ux430ux43dux438ux44f-ux445ux43eux441ux442ux430}

Если при установке виртуальной машины вы задали имяпользователя или имя
хоста, не удовлетворяющее соглашению обименовании, то вам необходимо
исправить это.

Запустите виртуальную машину и залогиньтесь.

Нажмите комбинацию Win+Enter для запуска терминала.

Запустите терминальный мультиплексор tmux: tmux

Переключитесь на роль супер-пользователя: sudo -i

Создайте пользователя (вместо username укажите ваш логин в дисплейном
классе): adduser -G wheel username

Задайте пароль для пользователя (вместо username укажите ваш логин в
дисплейном классе): passwd username

Установите имя хоста (вместо username укажите ваш логин в дисплейном
классе): hostnamectl set-hostname username

\subsubsection{Проверьте, что имя хоста установлено
верно:}\label{ux43fux440ux43eux432ux435ux440ux44cux442ux435-ux447ux442ux43e-ux438ux43cux44f-ux445ux43eux441ux442ux430-ux443ux441ux442ux430ux43dux43eux432ux43bux435ux43dux43e-ux432ux435ux440ux43dux43e}

hostnamectl (рис.~\ref{fig-006}).

\begin{figure}

\centering{

\includegraphics[width=0.7\linewidth,height=\textheight,keepaspectratio]{image/6.jpg}

}

\caption{\label{fig-006}Имя хоста}

\end{figure}%

\subsection{Установка программного обеспечения для создания
документации}\label{ux443ux441ux442ux430ux43dux43eux432ux43aux430-ux43fux440ux43eux433ux440ux430ux43cux43cux43dux43eux433ux43e-ux43eux431ux435ux441ux43fux435ux447ux435ux43dux438ux44f-ux434ux43bux44f-ux441ux43eux437ux434ux430ux43dux438ux44f-ux434ux43eux43aux443ux43cux435ux43dux442ux430ux446ux438ux438}

Нажмите комбинацию Win+Enter для запуска терминала.

Запустите терминальный мультиплексор tmux: tmux

Переключитесь на роль супер-пользователя: sudo -i

\subsection{Работа с языком разметки
Markdown}\label{ux440ux430ux431ux43eux442ux430-ux441-ux44fux437ux44bux43aux43eux43c-ux440ux430ux437ux43cux435ux442ux43aux438-markdown}

Средство pandoc для работы с языком разметки Markdown.

Установка с помощью менеджера пакетов: sudo dnf -y install pandoc

Для работы с перекрёстными ссылками мы используем пакет pandoc-crossref.
Пакет pandoc-crossref в стандартном репозитории отсутствует. Придётся
ставить вручную, скачав с сайта
https://github.com/lierdakil/pandoc-crossref.

При установке pandoc-crossref следует обращать внимание, для какой
версии pandoc он скомпилён. Лучше установить pandoc и pandoc-crossref
вручную.

Скачайте необходимую версию pandoc-crossref
(https://github.com/lierdakil/pandoc-crossref/releases). Посмотрите, для
какой версии откомпилён pandoc-crossref. Скачайте соответствующую версию
pandoc (https://github.com/jgm/pandoc/releases). Распакуйте архивы. Обе
программы собраны в виде статически-линкованных бинарных файлов.
Поместите их в каталог /usr/local/bin

\subsubsection{\texorpdfstring{(рис.~\ref{fig-007}).}{(рис.~).}}\label{ux440ux438ux441.-fig-007.}

\begin{figure}

\centering{

\includegraphics[width=0.7\linewidth,height=\textheight,keepaspectratio]{image/7.jpg}

}

\caption{\label{fig-007}Pandoc и pandoc-crossref}

\end{figure}%

\subsection{texlive}\label{texlive}

Установим дистрибутив TeXlive: sudo dnf -y install texlive-scheme-full

\#Домашнее задание

\subsubsection{\texorpdfstring{(рис.~\ref{fig-008}).}{(рис.~).}}\label{ux440ux438ux441.-fig-008.}

\begin{figure}

\centering{

\includegraphics[width=0.7\linewidth,height=\textheight,keepaspectratio]{image/8.jpg}

}

\caption{\label{fig-008}Домашнее задание(1)}

\end{figure}%

\subsubsection{\texorpdfstring{(рис.~\ref{fig-009}).}{(рис.~).}}\label{ux440ux438ux441.-fig-009.}

\begin{figure}

\centering{

\includegraphics[width=0.7\linewidth,height=\textheight,keepaspectratio]{image/9.jpg}

}

\caption{\label{fig-009}Домашнее задание(2)}

\end{figure}%

\section{Выводы}\label{ux432ux44bux432ux43eux434ux44b}

В результате выполнения данной лабораторной работы мы провели установку
системы Fedora Sway на виртуальную машину и её первоначальную настройку.




\end{document}
